% !TEX program = xelatex
\documentclass[10pt,aspectratio=169]{beamer}

\usetheme[
  lang=en,
  font=calibri,
  elevatedtitles=false,
  serifmath=true,
  tightlists=true,
  footertitle=true
]{OU}

% These options are passed to \usetheme[...]{OU}.
%
% lang=en|nl (default=en)
%   Selects English or Dutch labels and institute name.
%
% font=calibri|carlito|arial|none (default=arial)
%   Selects the main presentation font.
%   calibri : uses Calibri if available, with Carlito as fallback.
%   carlito : uses Carlito directly.
%   arial   : uses Arial if available, with Carlito as fallback.
%   none    : does not let the theme load or change fonts.
%
% elevatedtitles=true|false (default=false)
%   true  : moves frame titles closer to the top of the slide.
%   false : gives frame titles more vertical space.
%
% serifmath=true|false (default=true)
%   true  : keeps mathematical expressions in the usual LaTeX serif style.
%   false : uses the surrounding sans-serif style more strongly.
%
% tightlists=true|false (default=true)
%   true  : slightly reduces vertical spacing in standard itemize lists.
%   false : keeps Beamer's default list spacing.
%
% footertitle=true|false (default=true)
%   true  : shows the short title in the footer.
%   false : hides the short title in the footer.
%
% Short title used in the footer:
%   \title[Short footer title]{Full presentation title}
%
% Optional footer-title override:
%   \OUfootertitle{Different footer title}
%
% ------------------------------------------------------------
% OU per-frame options
% These options are passed to individual frames:
%
%   \begin{frame}[framecolor=OUAccent3]{Title}
%   ...
%   \end{frame}
%
% framecolor=<colour> (default=OUPrimary)
%   Changes the accent colour for one frame. The accent colour is used for
%   the logo mark and footer bar. It resets automatically on the next frame.
%
% Footline placement:
%   rightfootline : places the footer bar, page number, and optional short
%                   title on the right side of the frame (default).
%   leftfootline  : places the footer bar, page number, and optional short
%                   title on the left side of the frame.
%   nofootline    : removes the footer bar, page number, and optional short
%                   title. The logo layer is not affected.
%
% Logo placement:
%   leftlogo      : places the small OU logo mark on the left side of the
%                   title area (default). Additional frame logos are placed
%                   next to it automatically.
%   rightlogo     : places the small OU logo mark on the right side of the
%                   title area. Additional frame logos are placed next to it
%                   automatically.
%   nologo        : removes the regular-frame logo layer, including the small
%                   OU logo mark and any additional frame logos. The footline
%                   is not affected.
%
% Default frame layout:
%   rightfootline,leftlogo
%
% Examples:
%   \begin{frame}[leftfootline]{Title}
%   \begin{frame}[rightlogo]{Title}
%   \begin{frame}[leftfootline,rightlogo]{Title}
%   \begin{frame}[nofootline]{Title}
%   \begin{frame}[nologo]{Title}
%   \begin{frame}[nofootline,nologo]{Title}
%   \begin{frame}[framecolor=OUAccent3,leftfootline,rightlogo]{Title}
%
% Standard Beamer options such as t, plain, fragile, and noframenumbering
% can still be used together with the OU per-frame options.
%
% ------------------------------------------------------------
% OU additional logos
% ------------------------------------------------------------
% Additional logos can be added to the title page and to regular frames by
% giving only the logo filename. The optional argument is passed directly to
% \includegraphics, so standard graphic sizing options can be used.
%
% Add logos to the title page:
%   \OUaddtitlelogo{partner-logo.pdf}
%   \OUaddtitlelogo[height=0.65cm]{partner-logo.pdf}
%   \OUaddtitlelogo[width=2.0cm]{partner-logo.pdf}
%
% Add logos to regular frames:
%   \OUaddframelogo{partner-logo.pdf}
%   \OUaddframelogo[height=0.38cm]{partner-logo.pdf}
%   \OUaddframelogo[width=1.2cm]{partner-logo.pdf}
%
% Multiple logos can be added by repeating the commands:
%   \OUaddtitlelogo[height=0.65cm]{logo-a.pdf}
%   \OUaddtitlelogo[height=0.65cm]{logo-b.pdf}
%
%   \OUaddframelogo[height=0.38cm]{logo-a.pdf}
%   \OUaddframelogo[height=0.38cm]{logo-b.pdf}
%
% Title-page logos are placed automatically in the upper-right corner.
% Regular-frame logos are placed automatically next to the OU logo mark:
%   - with leftlogo, additional logos are placed to the right of the OU mark;
%   - with rightlogo, additional logos are placed to the left of the OU mark.
%
% Logo spacing can be adjusted without TikZ:
%   \setlength{\OUtitlelogogap}{0.6cm}
%   \setlength{\OUframelogogap}{0.45cm}
%
% Clear all additional logos:
%   \OUcleartitlelogos
%   \OUclearframelogos
%
% Advanced manual logo hooks are also available for custom TikZ placement:
%   \OUtitlelogos{<raw TikZ code>}
%   \OUframelogos{<raw TikZ code>}
%
% These advanced hooks replace the automatically generated additional-logo
% layer. They are intended only for layouts that cannot be expressed through
% the filename-based commands above.
%
% ------------------------------------------------------------
% OU colours and custom environments
% ------------------------------------------------------------
% Defined OU colours:
%   OUPrimary, OUAccent2, OUAccent3, OUwarmgrey, OUgreenish,
%   OUlightblue, OUblack, OUwhite
%
% Custom outline environment:
%   \begin{ououtline}
%     \item<1-> Main topic
%     \itemdetails<1-> Short explanation
%   \end{ououtline}
%
% Optional colour override:
%   \begin{ououtline}[OUAccent3] ... \end{ououtline}
% ------------------------------------------------------------

% ------------------------------------------------------------
% Common scientific-presentation packages
% ------------------------------------------------------------
\usepackage{amsmath,amssymb}
\usepackage{booktabs}
\usepackage{graphicx}
\usepackage{tikz}

% Helpful package to manage units and quantities consistently
\usepackage{siunitx}
\sisetup{
  per-mode=reciprocal,
  detect-all,
  number-mode=text
}

% ------------------------------------------------------------
% Title information
% ------------------------------------------------------------
\title[OU Beamer Template]{OU Beamer Template}
\subtitle{A short example presentation}
\author{Your Name}
\date{\today}

\begin{document}

% ------------------------------------------------------------
% Title page
% ------------------------------------------------------------
\begin{frame}[plain,noframenumbering]
  \titlepage
\end{frame}

% ------------------------------------------------------------
% Custom outline slide
% ------------------------------------------------------------
\section{Outline}
\begin{frame}[t]{Outline}
\setbeamercovered{dynamic}

\vspace{-0.5cm}
\begin{ououtline}
  \item Motivation
  \itemdetails Why the question matters and what is being investigated.
  \item Method
  \itemdetails A compact description of the model, data, or experiment.
  \item Results
  \itemdetails The main quantitative findings and their interpretation.
  \item Discussion
  \itemdetails Limitations, assumptions, and possible extensions.
\end{ououtline}
\end{frame}

% ------------------------------------------------------------
\section{Formatting}
\subsection{Text, figure, and table placement}
\begin{frame}[t]{Single-column text}
This is a normal text slide. Use short paragraphs and avoid overloading the slide with detailed prose. Inline mathematics such as \(E=mc^2\) and quantities such as \qty{3.5}{MW} are typeset consistently.

\medskip
The OU theme defines the title area, logo mark, and footer automatically. Most slides can therefore remain very simple in their source code.
\end{frame}

% ------------------------------------------------------------
\begin{frame}[t]{Columns and short lists}
\begin{columns}[T,onlytextwidth]
  \column{0.52\textwidth}
  \textbf{Left column}
  \begin{itemize}
    \item Use columns for comparisons.
    \item Keep lists short.
    \item Prefer readable slides over dense slides.
  \end{itemize}

  \column{0.42\textwidth}
  \textbf{Right column}

  \begin{tikzpicture}
    \draw[draw=OUPrimary,fill=OUwarmgrey,line width=0.5pt]
      (0,0) rectangle (4.8,2.8);
    \node[align=center] at (2.4,1.4) {Figure or diagram\\placeholder};
  \end{tikzpicture}
\end{columns}
\end{frame}

% ------------------------------------------------------------
\subsection{Equations and units}
\begin{frame}[t]{Equations and units}
A short derivation is often clearer than a dense list of equations. The example below uses \texttt{siunitx} for units.

%
\begin{equation}
\begin{aligned}
P &= \tau \omega \\
\tau &= F r \\
P &= F r \omega
\end{aligned}
\end{equation}
%

For example, with \(F=\qty{50}{kN}\), \(r=\qty{2}{m}\), and \(\omega=\qty{6}{rad.s^{-1}}\), the resulting power is \(P=\qty{600}{kW}\).
\end{frame}

% ------------------------------------------------------------
\begin{frame}[t]{Compact table}
\begin{table}
\centering
\small
\begin{tabular}{@{}lcc@{}}
\toprule
Case & Capacity & Comment \\
\midrule
Baseline & \qty{2.4}{GW} & reference case \\
Scenario A & \qty{2.8}{GW} & moderate increase \\
Scenario B & \qty{3.1}{GW} & strong increase \\
\bottomrule
\end{tabular}
\caption{A compact table using \texttt{booktabs}.}
\end{table}
\end{frame}

% ------------------------------------------------------------
\section{Beamer properties}
\subsection{Frame colour}
\begin{frame}[framecolor=OUAccent3]{Changing the frame colour}
This frame uses the per-frame option \texttt{framecolor=OUAccent3}. The accent colour returns to the default on the next frame.

\begin{block}{Useful for section changes}
Changing the frame colour can help mark exercises, interludes, case studies, or discussion slides without changing the whole presentation theme.
\end{block}
\end{frame}

% ------------------------------------------------------------
\begin{frame}[t]{Back to the default colour}
The accent colour is again \texttt{OUPrimary}, because \texttt{framecolor} is local to the frame.

\begin{exampleblock}{Example block}
Blocks can be used for definitions, assumptions, or short conclusions.
\end{exampleblock}
\end{frame}

% ------------------------------------------------------------
\subsection{Footline and logo placement}
\begin{frame}[t]{Default layout: \texttt{rightfootline,leftlogo}}
By default, the footer is placed on the right side of the frame and the OU logo mark is placed on the left side of the title area.

This corresponds to the implicit frame options \texttt{rightfootline,leftlogo}.
\end{frame}

% ------------------------------------------------------------
\begin{frame}[t,leftfootline]{Using \texttt{leftfootline}}
This frame uses the \texttt{leftfootline} option to move only the footer content to the left side of the frame.

The logo remains on the default left side, because footer placement and logo placement are controlled independently.
\end{frame}

% ------------------------------------------------------------
\begin{frame}[t,rightlogo]{Using \texttt{rightlogo}}
This frame uses the \texttt{rightlogo} option to move only the logo mark to the right side of the title area.

The footer remains on the default right side, because logo placement and footer placement are controlled independently.
\end{frame}

% ------------------------------------------------------------
\begin{frame}[t,leftfootline,rightlogo]{Combining \texttt{leftfootline} and \texttt{rightlogo}}
This frame uses \texttt{leftfootline,rightlogo}. The footer is placed on the left side, while the logo mark is placed on the right side.

This is useful when a figure or a diagram needs more space in one of the title-area corners.
\end{frame}

% ------------------------------------------------------------
\begin{frame}[t,nofootline]{Using \texttt{nofootline}}
This frame uses the \texttt{nofootline} option. The footer bar, page number, and optional short title are removed.

The logo mark remains visible, because \texttt{nofootline} does not affect the logo.
\end{frame}

% ------------------------------------------------------------
\begin{frame}[t,nologo]{Using \texttt{nologo}}
This frame uses the \texttt{nologo} option. The logo mark is removed from the title area.

The footer remains visible, because \texttt{nologo} does not affect the footline.
\end{frame}

% ------------------------------------------------------------
\begin{frame}[t,nofootline,nologo]{Using \texttt{nofootline} and \texttt{nologo}}
This frame combines \texttt{nofootline,nologo}. Both the footer and the logo mark are removed.

This can be useful for full-frame figures, diagrams, or interlude slides.
\end{frame}

% ------------------------------------------------------------
\begin{frame}[t,framecolor=OUAccent3,leftfootline,rightlogo]{Combining layout and colour}
This frame combines several OU per-frame options:

\medskip
\texttt{framecolor=OUAccent3,leftfootline,rightlogo}

\medskip
The accent colour, footer placement, and logo placement are all local to this frame.
\end{frame}

% ------------------------------------------------------------
\section{Remarks}
\begin{frame}[t]{Do you need a particular feature?}
  \begin{quote}
    If a particular feature is missing from the template, write to leonardo.rydin.gorjao@ou.nl.
  \end{quote}
\end{frame}

% ------------------------------------------------------------
\appendix
\begin{frame}[t]{Backup slide}
Backup slides can contain additional derivations, sensitivity analyses, or supporting figures.
\end{frame}

\end{document}